\chapter*{مقدمة هذه النسخة}
\addcontentsline{toc}{chapter}{مقدمة هذه النسخة}

هذه النسخة المحررة من دروس الشيخ \rawi{حسين بن عبدالرازق المصري} في قراءة «تفسير الطبري» أُعدَّت لتكون أقرب إلى صورة كتاب علمي تدريسي ينتفع به طالب العلم في المراجعة، والدرس، والتقييد، والرجوع السريع إلى المسائل.

ولأجل هذا المقصد جرى في العمل على هذه النسخة أمور تحريرية مقصودة، ينبغي التنبيه عليها بوضوح:
\begin{itemize}
    \item هذه المقدمة والخاتمة والملاحق ليست من كلام الشيخ، وإنما هي من عمل محرر هذه النسخة.
    \item الأصل المعتمد في المتن هو دروس الشيخ المفرغة في الملفات الخام، مع تصحيح أخطاء التفريغ الآلي، وحذف ما لا يحتاج إليه أسلوب الكتاب من لزمات التخاطب والتكرار الشفهي.
    \item في بعض المواضع وقع \textbf{تصرف تحريري} في العبارة: بتفكيك المعنى، أو جمع المتفرق، أو إعادة ترتيب الجمل، أو الإيضاح بلفظ أقرب إلى أسلوب التصنيف؛ مع المحافظة على المعنى المقصود في الجملة، لا سيما حيث كان السياق الشفهي يحتاج إلى إعادة سبك ليصلح للقراءة المتصلة.
    \item ما كان من زيادة تفسيرية أو تربوية أو منهجية من خارج لفظ الشيخ أُلحق في الغالب في صورة تنظيم للفقرات، أو عناوين، أو فوائد مستنبطة، أو ملاحق جامعة، ولم يقصد بذلك أن ينسب إلى الشيخ لفظاً لم يتلفظ به.
    \item لا تغني هذه النسخة عن الرجوع إلى الصوت الأصلي أو التفريغ الخام في المواضع التي يراد فيها التثبت من اللفظ أو تتبع مسار الاستدلال الشفهي كاملاً.
\end{itemize}

وقد ظهر في هذه السلسلة المباركة عدد من الخصائص التي تستحق التنبيه:
\begin{itemize}
    \item أنها لا تقتصر على حكاية ترجيحات \rawi{الطبري}، بل تكشف طريقته في النظر، والجمع، والنقد، والتوجيه.
    \item أنها تربط بين التفسير والفقه والأصول وعلوم القرآن والتزكية ربطاً حياً، فلا يبقى التفسير علماً معزولاً عن أثره العملي.
    \item أنها نافعة جداً في تعليم كيفية قراءة الكتب الكبار: كيف تجمع الأقوال، وكيف تفهم مأخذ العالم، وكيف تميز بين اختلاف التنوع والتضاد، وكيف تضبط السياق.
\end{itemize}

والرجاء أن تكون هذه النسخة خادمة لكلام الشيخ لا مزاحمة له، ميسرة لفهمه، جامعة لفوائده، باعثة على مزيد من العناية بتراث الأئمة المحققين، وعلى رأسهم الإمام \rawi{ابن جرير الطبري} رحمه الله.
